% Transcription — fonds Alexandre Grothendieck, shelfmark 8, batch 4 (pages 61-80).
% Established from the facsimile archives/batches/8/batch-04.pdf, cut from a
% copy of 8.pdf fetched through the project's own relay
% (https://grothendieck-relay.onrender.com/source/8.pdf) rather than from
% grothendieck.umontpellier.fr directly; per the skill transcribe-grothendieck.
% The batch file carries no cover sheet: PDF page k is the archivists' page
% 60+k, and their pencil numbers bottom-left agree throughout. The folder is
% eighty pages; this is the fourth and last batch (the others were
% transcribed separately by other passes).
%
% Pass: Opus 5.5 (claude-opus-5-5), 2026-09-26 — first pass, unchecked against the pages by a human.
%
% The folder sits in the inventory group "4-9" « Cristaux »; it has its own
% date in the catalogue, which \dating{} carries verbatim, with no group
% sentence.
%
% Pages skipped: 71 and 75, stencilled leaves « n° 397 » (pp. - 35 - and
% - 33 - of a text on measures, convolution and L^p spaces, Hölder's
% inequality), upside down, with nothing in his hand; 76, blank. Unrelated
% versos, skipped by the settled rule. (Batches 2-3 skipped a different
% stencilled text, « n° 382 », on C^r manifolds; none of its leaves occurs
% here.)
%
% Pages given as typescript, one French \note each with the hand-made
% interventions in full, the typed text not re-set: 61, 63, 65, 67, 69, 73.
% They are six leaves of a French typescript, numbered 35, 34, 33, 29, 28
% and 31 (61 and 73 scanned upside down), on "S-couples lisses" (Y, X) of
% codimension c, the Kummer sequence, the morphism epsilon_t of (3.9), a
% Théorème (3.10) giving canonical isomorphisms phi(Z, Y) for smooth pairs,
% its transitivity diagram (3.10.1) and its proof by induction on the
% codimension, and a base-change statement for (R^q j_*) g^* F proved via
% Hochschild-Serre. The Greek letters, the tildes over the isomorphism
% arrows, the accents and the cross-references (IX 3.2), (3.8), (VIII 5.5),
% (XII 6.5), (3.4) are inked or pencilled in by hand. Its page numbers
% collide with those of the proper-base-change typescript of batch 1
% (35, 37.1, 32, 33), so it is at least a different part, probably a
% different text; it is not identified here (no reference is looked up), and
% the hand of the corrections is not asserted.
%
% Pages transcribed from his hand: 62, 64, 66, 68, 70, 72, 74, 77-80.
% Pages summarised rather than set line by line: the sideways upper two
% thirds of 72 (scattered inscriptions under long diagonal strokes: H_B^*,
% H_Hdg, H_DR, H_cris, T_B, T_DR, T_Hdg), described in a \note with every
% legible inscription given.
%
% His pagination: none on the handwritten leaves; the typescript's own
% numbers are given in the notes.
%
% What the batch is about. Pp. 62-68 continue p. 60 (batch 3): the Galois
% group of the Tannakian category C of (filtered) F-crystals up to isogeny
% over V, with fibre functor Delta, T_p : C -> Mod_f(pi, Q_p) giving
% alpha : pi -> G(Q_l) (sic), G the algebraic envelope of alpha when T_p is
% fully faithful; the composite C -> Mod_f(pi, C) -> Mod_f(C) is bigraded
% (total degree and Hodge decomposition), whence i_1, i_2 : G_m,C => G_C; the
% Hodge fibre functor T_Hdg over K gives a torsor P under G_L with
% T_Hdg = P ^G (T_p (x) L) and j_1, j_2 : G_m,L => G'_L = ad(P); the
% Hodge-Tate decomposition T_p,C(M) = sum T_Hdg(M)^{i,j} (x) C(-i); a choice
% rho : C(1) = C gives a trivialisation rho~ in P(C), and a hom of
% G_m,C-schemes V(C(-1))^* -> P_C compatible with (i_1, i_2), (j_1, j_2);
% then the crystalline fibre functor T_cr with torsor Q under G_K, the
% filtration of T_cr (x) L with gr = T_Hdg, a torsor R under a unipotent
% U' inside Isom(Q_L, P), and Frobenius F : T_DR^sigma = T_DR, i.e.
% Q^sigma = Q. P. 70: the same structures for the category of (iso)motives
% over an open S of Spec(Z): Betti T_B with its group G over Q, l-adic
% realisations, pi = Gal(Qbar/Q) -> G(Q_p) unramified outside S, local
% Galois groups pi(p), pi(oo) = Z/2Z, i_1(p), i_2(p) and i_1(oo), i_2(oo),
% De Rham T_DR with torsor P and G' = ad(P), j_2, Frobenius F_p in G'(Q_p),
% F_oo in G'(R) of order 2, rho^(oo) in P(C), rho^(p)(j_2, lambda) in P(C^p).
% P. 72: the same list for a motive over S: a) module with integrable
% connection and filtration, b) Frobenius (crystals), c) real structure with
% F_oo, d) Betti local system, e) T_p, f) comparisons. P. 74: a list of
% conjectures (Hodge, Tate, Frobenius in De Rham cohomology, « double
% réseau », "Tate groups"). Pp. 77-80: « Complexe de cotg. des groupes de
% t.m. »: for G = D(M), M a lattice, the canonical resolution
% 0 -> R -> Z[M] -> M -> 0, the tori sequence 0 -> G -> G_m^M -> D(R) -> 0,
% Mazur-Roberts transitivity giving l_G = M (x)^L O_S (1.3), (1.4), (1.4'),
% the descent case f : S -> B_pi (1.6)-(1.8), then (2.1)-(2.2') omega_G,
% n_G = Tor_1(M, O_S), the class xi in Ext^2 and xi_Lambda in H^2(S, End(M)),
% its origin xi_0 in H^2(B_pi, End(M_0)), and « D'après Illusie » the
% obstruction to lifting the representation over Z/p^2 Z.
%
% Notation, aligned with batches 1-3: \mathscr{C} for his category (as on
% p. 60), C for the completed algebraic closure, \mathbb{C} for his
% double-stroked C (complex numbers, pp. 70-72), C^{p} / C^{(p)} as he writes
% them; \mathcal{P} for his round-bellied P (the torsor), Q and R plain;
% \boldsymbol{\pi} is not used here (he writes a plain pi for the Galois
% group on these pages, and it is kept plain); \ell^{G}_\bullet for his
% « l. » as in batch 1; \check{\mathbb{V}} for his double-stroked V with a
% hacek (p. 64).
%
% Hardest pages and readings to check first: p. 62, the subscript of
% G(Q_l) and « dans ce cas seulement »; p. 64, the subscript of Hom and the
% label on the boxed arrow; p. 66, the hatched block and its margin; p. 70,
% the title line and the C^{(p)} of the second i-line; p. 72 (sideways
% scribble); p. 74, the parenthesis read « (\ill{}, généralisée) »; p. 79,
% the lines after « Supposons p.ex. »; p. 69, the Roman numeral of
% « (\ill{} 4.5.) ».

\documentclass[11pt,a4paper]{article}
% Le préambule commun à toutes les transcriptions du fonds Grothendieck.
%
% Il tient en une idée : **l'incertitude se compose, elle ne se résout pas.**
% Une transcription qui lisse un mot douteux a détruit la seule chose pour
% laquelle on la faisait — un lecteur doit pouvoir savoir, à chaque endroit,
% ce qui était sur le feuillet et ce qui a été ajouté. Les six macros qui
% suivent sont donc l'appareil critique, et non de la décoration.
%
% Les mêmes commandes sont reconnues par `scripts/render.mjs`, qui produit la
% vue de lecture du site. Une macro ajoutée ici doit l'être aussi là-bas,
% faute de quoi le rendu échouera bruyamment — ce qui est voulu : un rendu
% silencieusement faux est pire qu'un rendu absent.

\ProvidesPackage{grothendieck}[2026/08/08 Transcriptions du fonds Grothendieck]

\usepackage{amsmath,amssymb,amsthm}
\usepackage{mathrsfs}
\usepackage{stmaryrd}
% Les diagrammes commutatifs sont partout dans ce fonds : c'est la langue
% même de la géométrie algébrique de Grothendieck, et une transcription qui
% les rendrait en prose ne transcrirait rien.
\usepackage{tikz-cd}
% Français par défaut — c'est la langue des feuillets — mais l'anglais est
% chargé aussi : la traduction et le résumé sont des documents anglais, et la
% typographie française (espaces avant les deux-points) y serait une faute.
% Ces documents basculent par \selectlanguage{english} en tête de corps.
\usepackage[english,french]{babel}
\usepackage{fontspec}
\usepackage{geometry}
\geometry{a4paper,margin=25mm,marginparwidth=28mm}
\usepackage{marginnote}
\usepackage{xcolor}
% ulem, not soul. `soul`'s \st and \ul are built on a character-by-character
% scan that XeLaTeX's font machinery breaks: the document fails with an
% undefined control sequence rather than degrading. ulem does the same two jobs
% and is Unicode-engine safe. `normalem` stops it hijacking \emph, which the
% transcriptions use for Grothendieck's own emphasis.
\usepackage[normalem]{ulem}
\usepackage{hyperref}
\hypersetup{colorlinks=true,linkcolor=black,urlcolor=black,pdfborder={0 0 0}}

% --- Métadonnées du lot -----------------------------------------------------
% Reprises telles quelles par le script de rendu, qui les affiche en tête de
% la vue de lecture. Elles fixent ce que le fichier prétend couvrir : sans
% elles, une transcription flotte sans point d'attache dans un fonds de seize
% mille feuillets.

\newcommand{\folder}[1]{\gdef\grothFolder{#1}}
\newcommand{\batch}[1]{\gdef\grothBatch{#1}}
\newcommand{\leaves}[2]{\gdef\grothFirst{#1}\gdef\grothLast{#2}}
\newcommand{\foldertitle}[1]{\gdef\grothFoldertitle{#1}}
\gdef\grothFolder{?}\gdef\grothBatch{?}\gdef\grothFirst{?}\gdef\grothLast{?}\gdef\grothFoldertitle{}

% --- Le résumé --------------------------------------------------------------
%
% Il ouvre la lecture modernisée, sous le titre « Résumé », et il est composé
% en petit. Ce n'est pas un ornement : c'est le seul passage du document qui
% ne soit pas des mathématiques, et un lecteur doit pouvoir le reconnaître
% avant de l'avoir lu — puis le sauter s'il connaît déjà le sujet, ou s'y
% arrêter s'il ne le connaît pas. Le filet qui le suit marque la reprise du
% corps.

\newenvironment{resume}
  {\small\setlength{\parskip}{0.45em}}
  {\par\medskip\hrule\medskip}

% --- Mots-clés ---------------------------------------------------------------
%
% En anglais, en clôture du résumé de la lecture modernisée : le vocabulaire
% moderne sous lequel chercher ce que les feuillets construisent. Le manifeste
% du site (`npm run manifest`) les extrait pour étiqueter la cote entière —
% cette ligne est la source unique des tags, il n'y a pas de fichier à côté.
\newcommand{\keywords}[1]{\par\smallskip\noindent
  {\footnotesize\textsc{Keywords} --- \textit{#1}}\par}

% --- L'appareil critique ----------------------------------------------------

% Le numéro de feuillet, en marge. C'est l'ancre qui permet de régler un
% désaccord sur une formule en regardant *un* feuillet plutôt qu'une chemise
% de six cents. Le site s'en sert aussi pour tourner les pages du fac-similé
% quand on fait défiler la transcription.
\newcommand{\leaf}[1]{%
  \marginnote{\footnotesize\color{gray!70}\textsf{#1}}%
  \label{leaf:#1}\ignorespaces}

% Illisible : on ne devine pas. Un mot inventé qui se lit comme les autres est
% le pire résultat possible.
\newcommand{\ill}{\textcolor{red!60!black}{[\,\ldots\,]}}

% Lecture douteuse : le mot est donné, et signalé comme incertain.
\newcommand{\uncertain}[1]{\uline{#1}}

% Ajout de l'éditeur — une lettre manquante, un mot sous-entendu.
\newcommand{\add}[1]{\textcolor{blue!50!black}{[#1]}}

% Biffé par Grothendieck lui-même. Ce qu'il a rayé fait partie du document :
% une rature dit souvent où la pensée a changé de direction.
\newcommand{\struck}[1]{\sout{#1}}

% Note du transcripteur, sur l'état matériel du feuillet.
\newcommand{\note}[1]{\par\smallskip{\small\color{gray!60!black}\textit{[#1]}}\par\smallskip}

% Note marginale de Grothendieck, distincte du corps du texte.
\newcommand{\marginal}[1]{\par\smallskip\noindent\hspace{1em}%
  {\small\color{gray!60!black}#1}\par\smallskip}

% --- Titre ------------------------------------------------------------------

\AtBeginDocument{%
  \begin{center}
    {\large\bfseries Fonds Alexandre Grothendieck --- cote n\textsuperscript{o}~\grothFolder}\\[2pt]
    {\itshape\grothFoldertitle}\\[4pt]
    {\small feuillets \grothFirst--\grothLast\ (lot \grothBatch)}\\[2pt]
    {\footnotesize\color{gray!60!black}Transcription établie d'après le fac-similé numérisé
     par l'Université de Montpellier.\\
     Les lectures incertaines sont soulignées, les illisibles notées [\,\ldots\,],
     les ajouts entre crochets.}
  \end{center}
  \vspace{1em}}

\endinput


\folder{8}
\batch{4}
\pages{61}{80}
\dating{[vers 1967- vers 1971]}
\watermark{Édition de démonstration}
\foldertitle{Cristaux (1970). Gribouillis et calculs divers : notes manuscrites (s.d.)}

\begin{document}

\page{61}

\note{feuillet de tapuscrit dactylographié, scanné tête-bêche, numéroté 35
en haut à droite, sans rapport direct avec les notes manuscrites voisines ;
il est gardé parce qu'il porte des interventions à la main, et le texte tapé
n'est pas recomposé. Le texte pose, pour une situation au-dessus de
$S' \to S$, $\psi(S,Z,Y,X) = \xi\,\eta\,\varphi(Z,Y)$, d'où $\psi$ de nature
locale et commutant aux changements de base ; si $(Z,X)$ est un couple lisse
de codimension $c > 1$, on trouve localement sur $X$ un triple $(Z,Y,X)$ avec
$0 < \operatorname{codim}(Y,X) < c$, et tout revient à démontrer que
$\psi(S,Z,Y,X)$ est indépendant de $Y$ (unicité en codimension $c$,
l'existence résultant de la descente) ; par une chasse au diagramme de
transitivité on se ramène au cas $\operatorname{codim}(Y,X) = 1$, avec deux
triples $(Z,Y_0,X)$, $(Z,Y_1,X)$, $Y_0$, $Y_1$ définis par $f_0 = 0$,
$f_1 = 0$ ; on pose $\bar S = \operatorname{Spec} \mathcal{O}_S[t]$,
$\bar X = X \times_S \bar S$, $\bar Z = Z \times_S \bar S$ et $\bar Y$ défini
par $(f_1 - f_0)t + f_0 = 0$ ; la page s'arrête sur « Alors $\bar Y/\bar S$
est lisse en les points au dessus des ». Interventions à la main : les
lettres $\psi$ (trois fois) et $\varphi$ tracées dans les blancs ; après le
signe $=$ de la formule, « $\xi$ » et un second signe, lu $\eta$ d'après la
page 65 ; le numéro 35. Un blanc laissé après « l'unicité de » n'est pas
rempli}

\page{62}

\note{suite de la page 60 (lot 3) : $\mathscr{C}$, $\Delta$, $T_p$, $\pi$,
$V$, $K$, $L$ comme là}

Pour simplifier, on supposera que $\mathscr{C}$ est de t.f.\ sur
$\mathbb{Q}_p$. De plus, on a un hom
\[
\boxed{\ \pi \xrightarrow{\ \alpha\ } G(\mathbb{Q}_\ell)\ }
\]
\note{l'indice se lit $\ell$ et non $p$ ; transcrit tel quel}
Si le foncteur \struck{\ill{}} $T_p$ (ou, ce qui revient au même,
\struck{\ill{}} le foncteur $\Delta$) est pl.\ fidèle, (dans
\uncertain{ce cas seulement}) $G$ est l'enveloppe algébrique de $\alpha$.

De plus, le composé
\[
\mathscr{C} \to \operatorname{Mod}_f(\pi, \mathbb{Q}_p) \xrightarrow{\ \otimes_{\mathbb{Q}_p} C\ } \operatorname{Mod}_f(\pi, C) \to \operatorname{Mod}_f(C)
\]
est un foncteur \emph{bigradué}, \struck{i.e.\ on a un} (gradué par degré
total, et par décomposition de Hodge) donc on a
\[
\boxed{\ \mathbb{G}_{m,C} \overset{i_1}{\underset{i_2}{\rightrightarrows}} G_C\ }
\]
avec
\[
i(\lambda) = i_1(\lambda)\, i_2(\lambda) = i(\lambda) : \mathbb{G}_{m,C} \to G_C
\]
\note{la première des deux lettres $i$ est surchargée ; lecture de l'égalité
douteuse}
l'hom \emph{central} définissant la graduation de $\mathscr{C}$.

De plus, on a un \struck{hom} foncteur fibre « bigradué de Hodge »
$T_{\mathrm{Hdg}}$ sur $K$, donc défini par
\[
\begin{cases}
\boxed{\ \mathcal{P},\ \text{torseur sous } G_L\ } \\
\bigl(T_{\mathrm{Hdg}} = \mathcal{P} \wedge^{G} (T_p \otimes_{\mathbb{Q}_p} L)\bigr)
\end{cases}
\]
\note{après la parenthèse, un mot biffé illisible}
et
\[
\begin{cases}
\mathbb{G}_{m,L} \overset{j_1}{\underset{j_2}{\rightrightarrows}} G'_L = \operatorname{ad}(\mathcal{P}) \\
j_1(\lambda)\, j_2(\lambda) = i(\lambda)
\end{cases}
\]
\note{$G'_L$ est écrit sur une lettre surchargée, avec un $\mathcal{P}$ en
exposant au-dessus}
avec un isom \struck{$T_{\mathrm{Hdg}}$}

\page{64}

\[
T_{p\,C}(M) = T_p(M) \otimes_{\mathbb{Q}_p} C \simeq \sum T_{\mathrm{Hdg}}(M)^{i,j} \otimes C(-i)
\]
Donc pour tout isom \struck{\ill{}}
\[
C(1) \overset{\rho}{\simeq} C
\]
on déduit un isom bigradué
\[
\tilde\rho : T_{p\,C}(M) \simeq T_{\mathrm{Hdg}\,C}
\]
i.e.\ une trivialisation
\[
\tilde\rho \in \mathcal{P}(C) = \operatorname{Hom}_{\operatorname{Spec} \ill{}}(\operatorname{Spec} C, \mathcal{P})
\]
\note{l'indice de « Hom » est « Spec » suivi d'une lettre surchargée}
D'autre part, si $\rho'$ est un autre isomorphisme, on a
\[
\rho' = u\rho, \qquad u \in C^*
\]
et on devra avoir
\[
\widetilde{u\rho} = i_1(u^{-1})
\]
Donc, \struck{on} \emph{on a} un hom de \struck{torseurs}
\emph{$C^*$-espaces \emph{à opérations} de $\pi$}
\[
\boxed{\ C(1)^* \longrightarrow \mathcal{P}(C)\ }
\]
\note{au-dessus de la flèche, deux mots lus « \uncertain{élément}
\ill{} » ; « à opérations » et « de $\pi$ » sont doublement soulignés}
quand on regarde $\mathcal{P}(C)$ comme espace où $C^*$ opère grâce à
$i_1^{-1}$ ; en d'autres termes, un hom de $\mathbb{G}_{m,C}$-schémas
\[
\check{\mathbb{V}}(C(-1))^* \longrightarrow \mathcal{P}_C,
\qquad
\check{\mathbb{V}}(C(-1))^* = \check{\mathbb{V}}(\mathbb{Z}_p(-1))^*_C
\]
\note{l'égalité est écrite verticalement sous la source ; une flèche courbe
va de $C(1)^*$ (dans l'encadré) à $\check{\mathbb{V}}(C(-1))^*$}
L'hom précédent doit être compatible avec $(j_1, j_2)$ d'une part,
$(i_1, i_2)$ d'autre part, i.e.\ $(j_1, j_2)$ résulte par \emph{descente}
de $i_1, i_2$, quand on choisit l'isom
\[
G_L^{\mathcal{P}} \otimes_L C \simeq G_C
\]
associé à un $\tilde\rho$ \ldots

\page{65}

\note{feuillet du même tapuscrit, numéroté 33. Il donne la suite spectrale
$(R^p v^!)(R^q u^!)F \Rightarrow (R^{p+q} w^!)F$ et l'isomorphisme
$\xi : (R^{2b} v^!)(R^{2a} u^!)\mu_n^{\otimes c} \xrightarrow{\sim}
(R^{2c} w^!)\mu_n^{\otimes c}$, puis le diagramme d'isomorphismes (3.10.1),
carré formé de $\varphi(Z,Y) : \mathbb{Z}/n \to (R^{2b}v^!)\mu_n^{\otimes b}$,
$\varphi(Z,X)$, $\eta = (R^{2b}v^!)(\varphi(Y,X) \otimes
\mu_n^{\otimes b})$ et $\xi$, « l'assertion de transitivité » étant sa
commutativité ; puis « Démonstration » : récurrence sur la codimension, à
partir de $c = 1$ ($c = 0$ trivial) ; le morphisme $\varepsilon_t$ de (3.9)
est de nature locale pour la topologie étale et commute aux changements de
base ; tout revient à démontrer que (3.9) est bijectif et indépendant du
choix de $t$, ce qui se vérifie fibre par fibre sur $Y$, et l'on peut
supposer $S = \operatorname{Spec} k$, $k$ séparablement clos. Interventions à
la main : $\xi$ (trois fois), $\varphi(Z,Y)$, $\varphi(Z,X)$, $\eta$ tracés
dans les blancs, et le tilde au-dessus de la flèche de $\xi$ ; « (3.8) » au
crayon après « $(R^q i^!)$ » ; accents ajoutés ; une lettre ajoutée au-dessus
de « cela » ; le numéro 33}

\page{66}

Il y a aussi un foncteur fibre $T_{\mathrm{cr}}$ \struck{au-dessus de}
dessus de $K$, à valeurs dans les vectoriels \struck{[gradués \&
filtrés]}. Donc on a un torseur
\[
\begin{cases}
\boxed{\ Q\ \text{torseur sous } G_K\ } \\
T_{\mathrm{cr}} \simeq Q \wedge^{G_K} (T_p)_K
\end{cases}
\]
\struck{$k_2 : \mathbb{G}_{m,K} \to G'' = G_K^{Q} = \operatorname{ad}(Q)$
définissant \ill{} \ill{} \ill{} de $G''$ \ldots\ $\in U''(K)$, $U''$ \ill{}
\ill{} \ill{} via $k_2$, d'où $k_2$ fournissant un torseur sous $U''$}
\note{le bloc, encadré et rattaché à l'accolade, est barré de hachures
obliques ; seuls quelques mots se lisent}

$T_{\mathrm{cr}} \otimes_K L$ a une filtration, et il y a un isom de
\struck{torseurs} foncteurs filtrés ($T_{\mathrm{Hdg}}$ filtré par $i_2$)
\[
\operatorname{gr}(T_{\mathrm{cr}} \otimes_K L) \simeq T_{\mathrm{Hdg}}
\]
\struck{des isom} un torseur $R$ sous \ldots\ correspondant à
\[
U' \subset G' = G_L^{\mathcal{P}}
\]
\marginal{\uncertain{gr.\ unipotent}}\note{la note marginale est reliée par
un trait à $U'$}
\struck{$\tilde k_2 : Q_L \xrightarrow{\sim} \mathcal{P}$ défini par $j_2$
satisfaisant $\tilde k_2 u = \tilde k_2 u \circ k_2$}
\note{ces deux lignes sont biffées d'un trait et de hachures ; « défini par
$j_2$ » reste hors des hachures}
\struck{i.e.\ on a un}
\[
\boxed{\ R \hookrightarrow \operatorname{Isom}(Q_L, \mathcal{P}) \qquad
(R \text{ torseur sous } U') \qquad G'' \supset U''\ }
\]
\note{dans « $\operatorname{Isom}(Q_L, \mathcal{P})$ » une lettre est
surchargée devant $Q_L$ ; à droite de l'encadré, un passage de quatre lignes
barré de croisillons se lit en partie : « il faudra expliquer \ill{} cela
correspond \ill{} $i_2$ \ill{} structures bigraduées »}

D'autre part, il y a sur le foncteur $T_{\mathrm{DR}}$ un Frobenius
\[
F : T_{\mathrm{DR}}^{\sigma} \xrightarrow{\ \sim\ } T_{\mathrm{DR}}
\]
\note{l'indice de $T$ dans la ligne « \ldots\ foncteur $T_{\mathrm{DR}}$ »
est surchargé}

\page{67}

\note{feuillet du même tapuscrit, numéroté 29. Il établit
$(u^*(R^1j_*)g^*F)_{\bar x} \approx ((R^1j_*)g^*F)_{\bar x} \approx
H^q(\tilde U, \tilde g^*\tilde F)$ ; puisque $\tilde U$ est fortement
$1$-acyclique pour $\mathbb{L}$, des isomorphismes (« suite spectrale de
Hochschild-Serre en bas dimensions ») $H^1(G, H^0(\tilde S, \tilde F))
\approx H^1(\tilde U, \tilde g^*\tilde F) \approx ((R^1j_*)g^*F)_{\bar x}$,
où $G = \varprojlim_n \mathbb{Z}/n\mathbb{Z} = G(\tilde U/\bar U)$ ;
$H^0(\tilde S, \tilde F)$ et $G$ commutant au changement de base
$\bar s \to S$, il en est de même de $(R^1j_*)g^*F$. Puis, pour $F$ de
$\mathbb{L}$-torsion, le théorème sous la forme
$u'^*(R^q i^!)g^*F \xrightarrow{\sim} (R^q i'^!)g'^*u^*F$ pour chaque $q$,
vrai pour $q = 2$, et pour $q$ quelconque si $F = \mathbb{Z}/n$ ($n$
premier aux caractéristiques résiduelles) par l'argument de récurrence de
la démonstration de (3.7) ; « on termine avec le lemme suivant ».
Interventions à la main : le tilde au-dessus de la flèche de la dernière
formule, et le « $!$ » de $i'^!$ repassé ; « (VIII 5.5) » au crayon dans
« Tenant compte de ( ) » ; accents ajoutés ; un long trait vertical au
crayon en marge gauche le long de la première moitié de la page}

\page{68}

qui correspond donc à un isom
\[
\boxed{\ Q^{\sigma} \simeq Q\ }
\]

\page{69}

\note{feuillet du même tapuscrit, numéroté 28. « Démonstration. La première
assertion veut naturellement dire que pour $S' \xrightarrow{u} S$
arbitraire, le morphisme $u'^*(R^qj_*)g^*F \to (R^qj'_*)g'^*u^*F$ est
bijectif dans les situations envisagées, où $u' = u \times_S X$, et la
dernière assertion équivaut à la première pour les faisceaux abéliens » ;
pour $q = 0$, les isomorphismes (3.2) $u'^*f^*F \simeq u'^*j_*g^*F$,
$f'^*u^*F \simeq j'_*g'^*u^*F$ donnent le résultat ; pour $q > 0$,
réduction à $S, X, Y$ de type fini sur $\operatorname{Spec}\mathbb{Z}$ et
$S' = \bar s$ point géométrique ; pour $F$ faisceau de ind-$\mathbb{L}$-groupes,
si la codimension de $(Y,X)$ n'est pas $1$, $(R^qj_*)g^* = (R^qj'_*)g'^* = 0$ ;
sinon on reprend les notations de la démonstration de (3.5). Interventions
à la main, au crayon : en marge gauche, contre « Démonstration », un
crochet et la remarque « démonstration peu naturelle », prolongée par un
trait vertical le long de tout le paragraphe ; les renvois « (\ill{} 4.5.) »
(chiffre romain barré dessus et dessous, lu « IX » ou « IV »), « (cf.\ XII
6.5) » et « (3.4) » écrits dans les blancs ; les tildes au-dessus des deux
flèches et le signe $\wr\wr$ vertical entre elles ; accents ajoutés ; un
trait vertical en marge devant « S,X,Y de type fini »}

\page{70}

Structures sur catégorie $\mathscr{C}$ des \uncertain{motifs} ($\tfrac12$
simples) sur un ouvert $S \subset \operatorname{Spec}(\mathbb{Z})$.
\note{« (iso) » est ajouté au-dessus de « motifs » ; « $\mathscr{C}$ » est
ajouté au-dessus de la ligne}

\marginal{$T_B$ \ldots\ $T_{\mathrm{DR}}$ \ldots\ « relations de $T_B$ et
$T_{\mathrm{DR}}$ »}\note{en marge gauche, trois accolades étiquetées
$T_B$, $T_{\mathrm{DR}}$ et, verticalement, « \uncertain{relations} de
$T_B$ et $T_{\mathrm{DR}}$ »}

$T_B$ : foncteur de Betti (sur $\mathbb{Q}$) relatif à : $\mathbb{C}$,
\quad $\boxed{\text{groupe } G \text{ sur } \mathbb{Q}}$

$T_\ell \simeq T_B \otimes_{\mathbb{Q}} \mathbb{Q}_\ell$, réalisation
$\ell$-adique (sur $\bar{\mathbb{Q}} \subset \mathbb{C}$)

\begin{tikzcd}[column sep=small, row sep=small]
\pi = \operatorname{Gal}(\bar{\mathbb{Q}}/\mathbb{Q}) \arrow[r] & G(\mathbb{Q}_p) \\
\pi(p) = \operatorname{Gal}(\bar{\mathbb{Q}}_p/\mathbb{Q}_p) \arrow[u] &
\end{tikzcd}

\note{à droite, « non ramifié dans $S - \lbrace p \rbrace$ » ; sous la
flèche horizontale, « $p \in S$ ($p$ fini ou $\infty$) » ; dans le cadre de
gauche, « $\pi(\infty) = \mathbb{Z}/2\mathbb{Z}$ » ; dans les deux premières
lignes, $\mathbb{Q}$ est écrit sur une lettre surchargée}

\[
\begin{aligned}
&i_1(\infty),\ i_2(\infty) : \mathbb{G}_{m,\mathbb{C}} \rightrightarrows G(\mathbb{C}) \\
&i_1(p),\ i_2(p) : \mathbb{G}_{m,C^{(p)}} \rightrightarrows G(C^{(p)})
\end{aligned}
\]
\marginal{\uncertain{résultats des autres structures} (cf.\ \ldots)}
\note{à droite : « $U^{(p)} \subset G_{C^{(p)}}$ » et « d'où
$U^{(\infty)} \subset G_{\mathbb{C}}$ (via $i_2^{(p)}$, $i_2^{(\infty)}$) » ;
les arguments $\infty$, $p$ des $i$ sont surchargés}

$T_{\mathrm{DR}}$ : foncteur de De Rham (sur $\mathbb{Q}$), groupe $G'$ sur
$\mathbb{Q}$, correspond à
$\boxed{\text{torseur } \mathcal{P} \text{ sous } G}$,
$G' = \operatorname{ad}(\mathcal{P}) = G^{\mathcal{P}}$.

\emph{filtré}, i.e.\ on a une famille d'hom
\[
j_2 : \mathbb{G}_{m,\mathbb{Q}} \to G'
\]
(classe des \ill{} définie mod.\ conjugaison par
$U' \subset G'$ ($\Leftrightarrow$ torseur $R$ sous $U'$)).
\note{« unipotent associé » est écrit sous $U' \subset G'$}

$\forall\, p \in S$, $p$ premier fini, \ldots
\[
\boxed{\begin{array}{l}
F_p \in G'(\mathbb{Q}_p) \\
F_\infty \in G'(\mathbb{R})
\end{array}}
\quad (\text{Frobenius})
\]
\note{sous $F_\infty$ : « d'ordre 2 »}

\emph{Élément} \struck{\ill{}}
$\boxed{\rho^{(\infty)} \in \mathcal{P}(\mathbb{C})}$ transformant la classe
des $j_2$ en celle des $i_2^{(\infty)}$.
\note{« Élément » est souligné, écrit sur un mot biffé ; « en celle des »
est ajouté au-dessus}

Pour $\forall\, j_2 \in U'$, et tout $\lambda \in C^{p}(1)^*$, un élément
\[
\boxed{\ \rho^{(p)}(j_2, \lambda) \in \mathcal{P}(C^{p})\ }
\]
transformant \struck{\ill{}} $j_2$ en $i_2$, i.e.
\[
\begin{cases}
C^{p}(1)^* \times U'(C^{p}) \longrightarrow \mathcal{P}(C^{p}) \\
\bigl[\ \mathbb{G}_{m,C^{p}} \xrightarrow{\ i_2^{(p)}\ } G'_{C^{p}}, \qquad G'_{C^{p}} \supset U'_{C^{p}}\ \bigr]
\end{cases}
\]
et compatibilités avec opérations de
\begin{enumerate}
\item[a)] $C^{p*} \times U'(C^{p})$
\item[b)] $\pi(p) = \operatorname{Gal}(\bar{\mathbb{Q}}_p/\mathbb{Q}_p)$
\end{enumerate}

\page{72}

\note{les deux tiers supérieurs de la page, écrits en travers (feuille
tournée d'un quart de tour), sont un griffonnage coupé de longs traits
obliques ; on y lit, sans ordre sûr : « $\mathbb{Q}$ » ;
« $H^*_B(X_{\mathbb{C}})$ --- $H^*_B(X) \otimes_{\mathbb{Q}} \mathbb{C}$ »,
avec « oui » au-dessus du trait ; « $H^*(X, \mathcal{O}_X)$ ---
$H^{\cdot\cdot}(X, \Omega^\cdot_X)$ » ; « $L$ \quad
$H_{\mathrm{Hdg}}(X)$ \quad $H_{\mathrm{DR}}$ \quad $L$ », deux mots entourés
au-dessus (« oui \ill{} » et « non ? ») ; « $H_{\mathrm{cris}}$ » ;
« $i_1$ » ; « $H_{\mathrm{DR}}(X) \otimes_{\mathbb{R}} \mathbb{C} \simeq
H(X_{\mathbb{C}}, \mathbb{R}) \otimes_{\mathbb{R}} \mathbb{C}$ » ;
« $H^*_{\mathrm{DR}}(X_L)$ --- $H^*_{\mathrm{DR}}(X_{\mathbb{C}})$ » ;
« $\sigma_1 \cdot \sigma_2 = f$ », « automorphisme » ; « $Q^\sigma \simeq Q$ » ;
« $X_{\mathbb{R}}$ --- $X_{\mathbb{C}}$ » ; « $G$ » ;
« $H^*(X(\mathbb{C}), \mathbb{C})$ » ; « $H^*(X$ » inachevé ; un (2)
entouré ; « $S \subset \operatorname{Spec}(\mathbb{Z})$ » ; « bigradué \ill{} » ;
« $\operatorname{Spec}\mathbb{Z}$ » ; enfin un petit tableau encadré à
accolade : « $T_B$ : $G_B$ ; $T_{\mathrm{DR}}$ : $G_{\mathrm{DR}}$ ;
$T_{\mathrm{Hdg}}$ »}

\begin{enumerate}
\item[a)] $M$ module à connexion intégrable \emph{sur} $S$, avec
\emph{filtration}.
\item[b)] $M(p)^{\sigma} \xrightarrow{\ F_p\ } M(p)$ Frobenius
\struck{\ill{}} (théorie des cristaux)
\item[c)] $M(\mathbb{R})$ filtré sur $S(\mathbb{R})$, muni d'un automorphisme
$F_\infty$ d'ordre $2$.\note{« analytique » est ajouté au-dessus de
« filtré », sur un mot biffé}
\item[d)] \struck{$B$}$T$ (Betti) système local sur
$S(\mathbb{C})^{\mathrm{an}}$ ($\Leftrightarrow \pi_1(S(\mathbb{C})^{\mathrm{an}}, s)
\to \operatorname{Aut}(T_s)$)
\item[e)] $T_p$ syst.\ local sur $S - S_p$ \qquad
$T_p \simeq T \otimes_{\mathbb{Z}} \mathbb{Z}_p$ sur $S_{\mathbb{C}}$
\item[$\mathrm{f}_\infty$)] $M_{\mathbb{C}}^{\mathrm{an}} \simeq T \otimes_{\mathbb{Z}} \mathbb{C}$
\item[$\mathrm{f}_p$)] --- --- --- ---
\end{enumerate}
\note{« $\mathrm{f}_\infty$ ) » est écrit sur un « f) » surchargé}

\page{73}

\note{feuillet du même tapuscrit, scanné tête-bêche, numéroté 31. Pour
$(Y,X)$ un $S$-couple lisse de codimension $1$, $Y$ défini par $t = 0$, la
suite exacte de Kummer $0 \to \mu_n \to \mathbb{G}_m \xrightarrow{n}
\mathbb{G}_m \to 0$ sur $U$ donne un morphisme « cobord »
$\delta : j_*\mathbb{G}_m \to R^1j_*\mu_n$, d'où une section $\delta t$ et,
$n$ annulant $R^1j_*\mu_n$, le morphisme
$(3.9)\ \varepsilon_t : (\mathbb{Z}/n)_Y \to (R^2i^!)(\mathbb{Z}/n)_X
\approx i^*(R^1j_*)(\mathbb{Z}/n)_U$, dont on vérifiera que c'est une
bijection canonique ; puis « Théorème (3.10) : Soit C la famille des couples
lisses de codimension constant. Il existe un fonction et un seul $\varphi$
donnant pour chaque $S$-couple lisse $(Y,X)$ de codimension $c$ un
isomorphisme (dit canonique) de faisceaux ». Interventions à la main :
« (IX 3.2) » dans le blanc après « Kummer » ; le $\Gamma$ de
$\Gamma(U, \mathcal{O}_U)$ corrigé sur un T ; la lettre $\varphi$ ; après
« codimension c », un signe d'insertion renvoyant à la marge gauche :
« et chaque $n$ premier aux caractéristiques résiduelles de $X$ » ; la
formule qui suit, « $(Y,X) : \mathbb{Z}/n \to \underline{H}^{2c}_Y(X, \mu$ »,
est barrée d'un trait ; accents ajoutés}

\page{74}

Conj.\ de Hodge (\ill{}, généralisée)

Conj.\ de Tate (\ill{}, généralisée)

Conj.\ des Frobenius dans Coh.\ de De Rham (à connexion)

\marginal{?}
Conj.\ du double réseau (revoir)

« Tate groups » (\ill{}, généralisé) sur schémas
$\begin{cases} \text{de t.f.\ sur } \mathbb{Z} \\ \text{de t.f.\ sur corps} \end{cases}$

« Tate grou\add{ps} » sur anneaux locaux d'inégales car.

\note{les trois indications de droite (« déc.\ de Hodge sur
$\mathbb{C}$ », « opération de Galois, conj.\ Frobenius », « Frobenius
cristallin ») sont écrites en colonne en face des trois premières
conjectures ; le « ? » est en marge devant « Conj.\ du double réseau »}

\page{77}

\section*{\uncertain{Complexe de cotg.}\ des groupes de t.m.}
\note{titre souligné}

1. \quad Soit $G = D(M)$, $M$ un réseau sur $S$ (\uncertain{$=$ gpe ct
tordu sur $S$}). Considérons la résolution plate canonique en deux crans
\[
(1.1) \qquad 0 \longrightarrow R \xrightarrow{\ v\ } \mathbb{Z}[M] \xrightarrow{\ u\ } M \longrightarrow 0 ,
\]
où $R = \operatorname{Ker} u$, $u([m]) = m$. On en déduit, par passage aux
duals, une suite exacte
\[
(1.2) \qquad 0 \to G \to T \to T' \to 0, \qquad T \simeq \mathbb{G}_m^{M}, \quad T' = D(R)
\]
\note{$T \simeq \mathbb{G}_m^M$ et $T' = D(R)$ sont écrits verticalement sous
$T$, $T'$}
résolvant $G$ par des tores. (c'est de dim.\ relative infinie \ldots)

Le triangle de transitivité de Mazur--Roberts donne un isomorphisme \struck{\ill{}}
\[
\ell^{G}_\bullet \simeq [\omega_{T'} \to \omega_T] \simeq [R \otimes_{\mathbb{Z}} \mathcal{O}_S \to \mathbb{Z}[M] \otimes_{\mathbb{Z}} \mathcal{O}_S]
\]
i.e.
\[
(1.3) \qquad \boxed{\ \ell^{G}_\bullet \simeq \underline{M} \overset{L}{\otimes}_{\mathbb{Z}} \mathcal{O}_S\ }
\]
\note{après l'encadré, « etc » biffé}
Si \struck{\ill{}} $\mathcal{O}_S$ est une $\Lambda$-Algèbre (p.\,ex.\
$\Lambda = \mathbb{Q}$, ou $\Lambda = \mathbb{Z}/n\mathbb{Z}$)
\struck{, $n\mathcal{O}_S = 0$, i.e.\ $\mathcal{O}_S$ une
$\mathbb{Z}/n\mathbb{Z}$-Algèbre,} on peut écrire aussi
\[
(1.4) \qquad \boxed{\ \ell^{G}_\bullet \simeq \bigl(M \overset{L}{\otimes}_{\mathbb{Z}} \Lambda\bigr) \overset{L}{\otimes}_{\Lambda} \mathcal{O}_S\ }
\]
\note{les deux $\Lambda$ sont écrits sur des lettres biffées}
et si

\page{78}

\struck{$n = p$ (p.\,ex.\ $\Lambda = \mathbb{F}_p$) \ill{} \ill{}}
\note{au-dessus de la ligne biffée : « si $\Lambda$ est un corps, »}
p.\,ex.\ si $\Lambda$ est un corps, le $L$ sur le dernier $\otimes$ est
inutile, on peut écrire
\[
(1.4') \qquad \ell^{G}_\bullet \simeq \bigl(M \overset{L}{\otimes}_{\mathbb{Z}} \mathbb{F}_p\bigr) \otimes_{\mathbb{F}_p} \mathcal{O}_S .
\]

1.5\note{« 1.5 » est écrit sur un « 2. »} Supposons qu'on ait un
revêtement principal $S'$ de $S$ de groupe \struck{ordinaire} $\pi$ discret
ou profini, et que $M$ provienne par descente d'une \struck{revêtement}
représentation de $\pi$ dans un groupe abélien $M_0$ (qu'on suppose
continue pour $\pi$ profini) ; p.\,ex.\ si $M$ est isotrivial (p.\,ex.\ à
fibres de t.\ finie et $S$ \struck{connexe} normal) et si $S$ connexe,
\ill{} \ill{} dans la situation dite, avec $S'$ le rev.\ universel de
$S$ \ldots\ On a donc un hom.\ de topos
\[
(1.6) \qquad S \xrightarrow{\ f\ } (B_\pi)
\]
\note{sous $S$ : « (top.\ étale, ou fppf, ou fpqc) » ; après $S$ et après $B_\pi$, une lettre biffée illisible}
avec
\[
(1.7) \qquad M = f^{-1}(M_0)
\]
\note{l'exposant $-1$ est écrit sur une étoile}
On aura alors, si on annèle $S$ et $B_\pi$ par l'anneau constant $\Lambda$,
et envisage $f$ en conséquence
\[
M \overset{L}{\otimes}_{\mathbb{Z}} \Lambda = f^*\bigl(M_0 \overset{L}{\otimes}_{\mathbb{Z}} \Lambda\bigr)
\]
\note{numéro de formule surchargé, illisible}
et par suite on peut écrire
\[
(1.8) \qquad \ell^{G}_\bullet \simeq f^*\bigl(M_0 \overset{L}{\otimes}_{\mathbb{Z}} \Lambda\bigr) \overset{L}{\otimes}_{\Lambda} \mathcal{O}_S ,
\]
où le $L$ sur $\otimes_\Lambda$ est inutile si $\Lambda$ est un
corps \ldots

\page{79}

2. Notons que \struck{si $\mathcal{O}_S$ est de car.\ \ill{}} l'on tire de
(1.3)
\[
(2.1) \qquad
\begin{cases}
\omega_G = \underline{M} \otimes_{\mathbb{Z}} \mathcal{O}_S \\
\underline{n}_G = \underline{\operatorname{Tor}}^{\mathbb{Z}}_1(M, \mathcal{O}_S)
\end{cases}
\]
ce qui s'écrit, si $S$ est de car.\ $p > 0$
\[
(2.2) \qquad
\begin{cases}
\omega_G = (M \otimes_{\mathbb{Z}} \mathbb{F}_p) \otimes_{\mathbb{F}_p} \mathcal{O}_S = M_p \otimes_{\mathbb{F}_p} \mathcal{O}_S \\
\underline{n}_G \simeq {}_pM \otimes_{\mathbb{F}_p} \mathcal{O}_S
\end{cases}
\]
et si $M$ est annulé par $p$, i.e.\ $G$ annulé par $p$, cela devient
\[
(2.2') \qquad \omega_G \simeq M \otimes_{\mathbb{F}_p} \mathcal{O}_S, \qquad \underline{n}_G \simeq M \otimes_{\mathbb{F}_p} \mathcal{O}_S
\]
On a d'autre part une classe canonique
\[
\xi \in \operatorname{Ext}^2_{\mathcal{O}_S}\bigl(M \otimes_{\mathbb{Z}} \mathcal{O}_S, \operatorname{Tor}^{\mathbb{Z}}_1(M, \mathcal{O}_S)\bigr)
\]
qu'on peut interpréter, en présence d'un $\Lambda$, comme provenant d'un
\[
\xi_\Lambda \in \operatorname{Ext}^2_\Lambda\bigl(M \otimes_{\mathbb{Z}} \Lambda, \operatorname{Tor}^{\mathbb{Z}}_1(M, \Lambda)\bigr) ;
\]
Supposons p.\,ex.\ que $\Lambda = \mathbb{Z}/n\mathbb{Z}$, et que $M$
\struck{$M \otimes \mathbb{Z}$} soit un $\Lambda$-module loc.\ libre de t.f.
Alors \ill{} \ill{} \ill{} de $M \otimes_{\mathbb{Z}} \Lambda \simeq M$ et de
$\operatorname{Tor}^{\mathbb{Z}}_1(M, \Lambda) \simeq {}_nM \simeq M$, et on a
donc
\[
\xi_\Lambda \in H^2(S, \operatorname{End}(M)) ,
\]
et \struck{si $\mathcal{O}_S$ est plat sur $\Lambda$ (p.\,ex.\ \ill{})}
\[
\xi \in H^2(S, \operatorname{End}(E)) \qquad (E = M \otimes_\Lambda \mathcal{O}_S)
\]
et $\xi$ se déduit de $\xi_\Lambda$ par l'hom.\ can.

\page{80}

\[
H^2(S, \operatorname{End}(M)) \to H^2(S, \operatorname{End}(M) \otimes_\Lambda \mathcal{O}_S) \simeq H^2(S, \operatorname{End}(E))
\]
Enfin, sous les conditions de 1.5, $M_0$ étant aussi plat sur $\Lambda$, on
aura que $\xi_\Lambda$ est l'image d'un élément
\[
\xi_0 \in H^2(B_\pi, \operatorname{End}(M_0)) .
\]
(NB Pour calculer $\xi_0$, par fonctorialité en $\pi$, on est ramené
aussitôt au cas « universel » $\pi = \operatorname{Aut}(M_0)$, et on peut
évidemment supposer $M_0 = \Lambda^d = (\mathbb{Z}/n\mathbb{Z})^d$, donc
$\pi = \mathrm{Gl}(d, \mathbb{Z}/n\mathbb{Z})$ \ldots) D'après
\emph{Illusie}, c'est l'obstruction à remonter la représentation de $\pi$
dans $M_0$ en une représentation dans un module libre sur
$\mathbb{Z}/p^2\mathbb{Z}$. Lorsque $n = p$ (premier), elle est évidemment
nulle si $M_0$ est de rang $\le 1$, mais sans doute pas nulle si $M_0$ de
rang $\ge 2$, dans le cas universel $\pi = \mathrm{Gl}(d, \mathbb{F}_p)$.
\note{« Lorsque $n = p$ (premier) » est ajouté au-dessus de « Elle », et
relié à la ligne}

Notons que si le rev.\ $S'$ de $S$ est le rev.\ universel, alors
\[
H^2(B_\pi, -) \longrightarrow H^2(S, -)
\]
est \emph{injectif}, donc on trouve un élément \ill{} une
\struck{donnée} $\xi_\Lambda \in H^2(S, \operatorname{End}(M))$.

\end{document}
